\documentclass[10pt,a4paper]{amsart}
\usepackage[margin=19mm]{geometry}
\usepackage{amsmath,amssymb,mathtools}
\usepackage[scr=rsfs,cal=euler]{mathalfa}
\usepackage{xfrac}
\usepackage[unicode,hidelinks,bookmarks=false]{hyperref}
\newcommand{\ie}{i.e.\ }
\newcommand{\cf}{cf.\ }
\newcommand{\resp}{respectively\ }
\newcommand{\Subsection}{\subsection}
\providecommand{\nobreakdash}{\nobreak\hbox{-}}
\setlength{\emergencystretch}{2em}
\multlinegap0pt
\usepackage{amssymb}
\usepackage{bm}
\usepackage{mathtools}
\usepackage{enumerate}
\usepackage{algorithm}\usepackage{algpseudocode}
\usepackage{xcolor}
\usepackage[normalem]{ulem}
\usepackage{tikz}
\usepackage{tcolorbox}
\newtheorem{lem}{Lemma}
\newtheorem{thm}{Theorem}
\def\C{{\mathbb C}}
\def\EE{{\mathbb E}}
\def\PP{{\mathbb P}}
\def\PR{\mathbb{P}}
\title{Maximal Minimal Spacing for Random Points}
\author{Fabio Deelan Cunden, Noemi Cuppone, Giovanni Gramegna, Pierpaolo Vivo}
\date{Source: arXiv:2606.04837v2 (CC BY 4.0)}
\begin{document}
\maketitle

\noindent\emph{Source and licence.} Maximal Minimal Spacing for Random Points. Version arXiv:2606.04837v2 (CC BY 4.0), distributed by arXiv under the Creative Commons Attribution 4.0 International licence, http://creativecommons.org/licenses/by/4.0/, as recorded in the arXiv licence metadata for this exact version and shown on the abstract page https://arxiv.org/abs/2606.04837v2. Full licence notice: \texttt{licenses/arxiv-2606.04837-CC-BY-4.0.txt}.\par\smallskip

\noindent\textit{Editorial scope.} Counted unit: the article's thm thm:LD (source0.tex lines 532--562). The article's complete original proof of it is delivered with it (source0.tex lines 1166--1320, one \texttt{proof} environment of 155 lines, which the article places away from the statement). No mathematics is written, repaired or re-derived: the statement and the proof are the article's own lines, in the article's own notation, and no retained line is substituted or re-flowed.  Retained uncounted as the source-local context the proof needs: lines 410--413; lines 443--452; lines 455--470; lines 516--522; lines 528--530.  One declared adaptation: exactly 1 ancillary strings inside the retained spans are omitted (image-inclusion commands and, where the source repeats a label, the repeated label definitions) and no other line differs from the source; the figure environments, with their placement, captions and labels, are kept, and the package records the omitted strings in fidelity receipts bound to the affected bodies.  No retained line contains a citation, so the wrapper carries no bibliography and no bibliography entry.  The retained lines contain the article's own diagram code, which the wrapper typesets with the standard tikz packages; no image file is delivered and the package declares no asset dependency.  Editorial formatting only, in addition: an AMS \texttt{amsart} preamble that restates the article's own declarations and macros used by the retained lines, namely the environments \texttt{lem}, \texttt{thm} on separate counters, together with the standard packages those lines need.  The retained environments print in this document as Lemma 1, Theorem 1 in order of appearance, whereas the article prints the same items with the numbering of its own full document.  The complete native source of the article as distributed in the arXiv e-print for this exact version is bundled unchanged as \texttt{sources/202078/source0.tex}.
\begin{lem}
\label{lem:equiv}
   $ S^*\geq s \iff K_N(s)\geq M \iff \tau_M(s)\leq N$ .
\end{lem}

\begin{thm}[Distribution-free formula for the max-min spacing]
\label{thm:main}
For all $1\leq M\leq N$:
\begin{equation}\label{eq:main}
\PR(S^* \geq s)
   = [z^N]\frac{p_s(z)^M}{1-z}.
%    = [z^N](1-z)^{-1}p_s(z)^M.
\end{equation}

\end{thm}

\begin{lem}[Telescopic formula] 
\label{lem:telescopic} The distribution of the first-passage time $\tau_1(s)$ is 
\begin{equation}
  \PR(\tau_1(s)=k)= \PR(P_{k-1}< s)-\PR(P_{k}< s).
  \label{eq:PL}
\end{equation}
In particular,
\begin{equation}
\begin{aligned}
       p_s(z)
        &=z-(1-z)\sum_{k\geq1}\PR(P_{k}< s)z^{k}.
    \label{eq:pgf-telescopic} 
\end{aligned}
\end{equation}
 Moreover, if $\PP(T_1=0)=0$, then $p_s(z)$ is analytic in the whole complex plane $z\in\C$.
\end{lem}

\begin{equation}\label{eq:TiltedMeanVar}
\mathbb{E}[\tau_1(s,z)]
= \frac{z\;p_s'(z)}{p_s(z)},
\qquad
\operatorname{Var}[\tau_1(s,z)]
 = z\frac{d}{dz}\mathbb{E}[\tau_1(s,z)].  
\end{equation}

\begin{equation}\label{eq:typical}
\alpha \mathbb{E}[\tau_1({s^*},1)]=1.
\end{equation}

\begin{thm}[Large deviations]
\label{thm:LD}
Let $(T_j)_{j\geq1}$ be i.i.d. strictly positive random variables, and 
let $\alpha\in (0,1)$. Assume that, for each value of $s$, 
the saddle-point equation
\begin{equation}
\label{eq:assump1}
    \alpha \EE[\tau_1(s,z)]=1
\end{equation}
admits a positive solution $z(s)>0$, and that
\begin{equation}
\label{eq:assump2}
    \operatorname{Var}(\tau_1(s,z(s)))>0.
\end{equation}
Let $s^*$ be defined by~\eqref{eq:typical}. Suppose that $N\to\infty$ with $M=\lfloor \alpha N\rfloor$. Then, for $s>s^*$,
\begin{equation}\label{eq:right-box}
 \PP(S^*\ge s) =
  \frac{e^{-N\psi(s)}}{(1{-}z(s))\,
  \sqrt{2\pi \alpha N\sigma^2(s)}}\left[1+o(1)\right],
\end{equation}
while for $s<s^*$,
\begin{equation}\label{eq:left-box}
\PP(S^*<s) =
  \frac{e^{-N\psi(s)}}{(z(s){-}1)\,
  \sqrt{2\pi \alpha N\sigma^2(s)}}\left[1+o(1)\right].
\end{equation}
Here, 
    \begin{equation}\label{eq:psi}
    \psi(s)=\ln z(s)-\alpha\log p_s(z(s)),\qquad\sigma^2(s) =\operatorname{Var}[\tau_1(s,z(s))]. 
    \end{equation}
\end{thm}

\begin{proof}[{Proof of Theorem~\ref{thm:LD}}] 


{By Theorem~\ref{thm:main}} and Cauchy integral formula,
\begin{equation}\label{eq:P}
\PR(S^*\ge s)
  = [z^N] \frac{p_s(z)^M}{(1-z)}=\frac{1}{2\pi \mathrm{i}}\oint_{\gamma}
  \frac{p_s(z)^M}{(1-z)\,z^{N+1}}\,dz= \frac{1}{2\pi \mathrm{i}}\oint_{\gamma}
  \frac{e^{Ng(z)}}{(1{-}z)\,z}\,dz ,
\end{equation}
where $\gamma$ is a positively oriented simple contour  entirely contained in the unit disk and enclosing $z=0$, and 
\begin{equation}
    g(z)=\alpha\log p_s(z)-\log z.
\end{equation} 
By Lemma~\ref{lem:telescopic}, the integrand in~\eqref{eq:P} is meromorphic with poles at $z=0$ and $z=1$, only. 
Choose $\gamma$ to be a circle $C_r$ of radius $r$ centred at $0$. 
Since  $p_s(z)$ is a power series with positive coefficients, 
\begin{equation}
    |p_s(r e^{\mathrm{i} \theta})|=\left|\sum_{k\geq1} \mathbb{P}(\tau_1(s)=k){r^k e^{\mathrm{i}k\theta}}\right|\leq\sum_{k\geq1} \mathbb{P}(\tau_1(s)=k){r^k} =p_s(r),
\end{equation}
therefore
$\operatorname{Re}g(r{e^{\mathrm{i}\theta}})
$ has a global maximum at $\theta=0$.
For all $s>0$, choose $r=z(s)$, where $z(s)$ is a real and positive solution of $g'(z)=0$, i.e. $ \mu(z,s)=1/\alpha$, where $$\mu(z,s)=z \frac{p'_s(z)}{p_s(z)}=\EE [\tau_1(s,z)] .$$

The meaning of this saddle-point equation is the following.  The event
$\{S^*\geq s\}$ is equivalent, by Lemma~\ref{lem:equiv}, to completing $M$
renewal cycles by time $N$.  Hence, if $M/N=\alpha$, the
typical length of one cycle is $N/M=1/\alpha$.  The tilted law
$\tau_1(s,z)$ is precisely the exponential change of measure under which the
first-passage time has mean
\begin{equation}
    \mu(z,s)=\EE[\tau_1(s,z)] .
\end{equation}
Thus the saddle-point condition
\begin{equation}
    \mu(z,s)=\frac{1}{\alpha}
\end{equation}
selects the value of the tilt for which a typical renewal cycle has the length
required by the constraint.  In particular, $z<1$ biases the renewal process
towards shorter cycles, while $z>1$ biases it towards longer cycles.

By assumption, for each $s>0$ there exists a positive
solution $z(s)>0$ of the saddle-point equation
\begin{equation}
    \mu(z,s)=\frac{1}{\alpha},
    \qquad
    \mu(z,s):=
    z\frac{p_s'(z)}{p_s(z)}
    =
    \EE[\tau_1(s,z)] .
\end{equation}
Moreover, from \eqref{eq:TiltedMeanVar},
\begin{equation}
    \frac{d}{dz}\mu(z,s)
    =
    \frac{1}{z}\operatorname{Var}[\tau_1(s,z)].
\end{equation}
Hence, whenever $\operatorname{Var}(\tau_1(s,z))>0$, such a solution is unique.

Computing the second derivative, we have \begin{equation}
\label{eq:var_tauz}
    g''(z(s))=\alpha \frac{\operatorname{Var}[\tau_1(s,z(s))]}{z(s)^2}>0 .
\end{equation} 
 
On the circle $z=z(s)e^{\mathrm{i}\theta}$, the saddle equation gives
\begin{equation}
    g(z(s)e^{\mathrm{i}\theta})
    =
    g(z(s))
    -
    \frac{\alpha\sigma^2(s)}{2}\theta^2
    +
    O(\theta^3),
    \qquad \theta\to0 ,
\end{equation}
where $\sigma^2(s)=\operatorname{Var}(\tau_1(s,z(s)))$.
Thus $\theta=0$ is a nondegenerate maximum of
$\operatorname{Re}g(z(s)e^{\mathrm{i}\theta})$ along the circular contour.



Let
\begin{equation}
    \label{eq:rate_psi_g}
\psi(s)=-g(z(s))=-\operatorname{Re}g(z(s)).
\end{equation}
(Since $z(s)>0$, the quantity $g(z(s))$ is real.) 
Finally, the steepest descent direction for $\operatorname{Re}g(z)$ is parallel to the imaginary axis, and it is therefore aligned with the tangent to  $C_{z(s)}$ at $z(s)$. We are now prepared to apply the steepest descent method. 
\begin{figure}

\caption{Contours of integration appearing in the calculation of $\PR(S^*\geq s)$}
\label{fig:pic}
\end{figure}
By definition, $s^*$ is the solution of $\mu(1,s^*)=1/\alpha$. Hence $z(s^*)=1$ and, by monotonicity, %, $z(s)<1$ for $s>s^*$, while $z(s)>1$ for $s<s^*$. 
    \begin{equation}\label{eq:wts}
    \begin{cases}
        &s<s^* \Rightarrow z(s)>1\\
        &s=s^* \Rightarrow z(s)=1\\
        &s>s^*\Rightarrow z(s)<1 .
    \end{cases}
\end{equation}
Therefore, if $s>s^*$, the contour $C_{z(s)}$ is contained in the unit disk where the integrand $\frac{e^{Ng(z)}}{(1-z)z}$ has a single pole at $z=0$ (see Figure~\ref{fig:pic}). Therefore,
\begin{equation}
\begin{aligned}
\PP(S^*\geq s)
&=
\frac{1}{2\pi}
\int_{-\pi}^{\pi}
\frac{
\exp\{N g(z(s)e^{\mathrm{i}\theta})\}
}
{1-z(s)e^{\mathrm{i}\theta}}
\,d\theta\\
&=
\frac{e^{-N\psi(s)}}{
(1-z(s))\sqrt{2\pi\alpha N\sigma^2(s)}
}
\left[1+o(1)\right],
\end{aligned}
\end{equation}
which yields~\eqref{eq:right-box}.

For $s<s^*$, the contour deformation from $\gamma$ to $C_{z(s)}$ crosses the pole at $z=1$. Therefore,
\begin{equation}
 \PR(S^*\ge s)
 =\frac{1}{2\pi \mathrm{i}}\oint_{\gamma}
  \frac{e^{Ng(z)}}{(1{-}z)\,z}\,dz= \frac{1}{2\pi \mathrm{i}}\oint_{C_{z(s)}}
  \frac{e^{Ng(z)}}{(1{-}z)\,z}\,dz-\operatorname{Res}_{z=1}\left(\frac{e^{Ng(z)}}{(1{-}z)\,z}\right) . 
\end{equation}
 The residue is:
 \begin{equation}
\operatorname{Res}_{z=1}\left(\frac{e^{Ng(z)}}{(1{-}z)\,z}\right)=\lim_{z\to 1}(z-1)\frac{p_s(z)^M}{(1-z)z^{N+1}}=-p_s(1)^M=-1.
 \end{equation}
 Therefore,
\begin{equation}
\begin{aligned}
\PP(S^*<s)=1- \PR(S^*\geq s)
&=
-\frac{1}{2\pi}
\int_{-\pi}^{\pi}
\frac{
\exp\{N g(z(s)e^{\mathrm{i}\theta})\}
}
{1-z(s)e^{\mathrm{i}\theta}}
\,d\theta  \\
&=
\frac{e^{-N\psi(s)}}{
(z(s)-1)\sqrt{2\pi\alpha N\sigma^2(s)}
}
\left[1+o(1)\right],
\end{aligned}
\end{equation}
which matches~\eqref{eq:left-box}.
\end{proof}

\end{document}
